\documentclass[tikz,border=10pt]{standalone}
\usetikzlibrary{intersections, calc}

\begin{document}
\begin{tikzpicture}[scale=1.2]
    % Define Center and Radius
    \coordinate (O) at (0,0);
    \def\R{5}
    
    % Define Points on Circle
    \coordinate (B) at (-\R, 0);
    \coordinate (C) at (\R, 0);
    \coordinate (H) at (0, -\R);
    
    % A is near top, slightly right
    \coordinate (A) at (0.8, 4.93); 
    % D is top left
    \coordinate (D) at (-4.2, 2.7);
    
    % Draw Circle
    \draw (O) circle (\R);
    
    % Draw Chords and Lines
    \draw[name path=chord_AB] (A) -- (B);
    \draw[name path=chord_AC] (A) -- (C);
    \draw[name path=chord_CD] (C) -- (D);
    \draw (A) -- (D);
    \draw (B) -- (D);
    \draw (B) -- (C); % Diameter
    \draw (B) -- (H);
    \draw (C) -- (H);
    
    % Intersection E of AB and CD
    \path[name intersections={of=chord_AB and chord_CD, by=E}];
    
    % Point G: Projection of H onto CD (perpendicular foot)
    \coordinate (G) at ($(C)!(H)!(D)$);
    
    % Line H-G-F. F is intersection of line HG extended with AC.
    \path[name path=line_HG] (H) -- ($(H)!2!(G)$);
    \path[name intersections={of=line_HG and chord_AC, by=F}];
    
    % Draw line segment HF (passing through G)
    \draw (H) -- (F);
    
    % Draw segment OG
    \draw (O) -- (G);
    
    % Right angle mark at G (between CD and HF)
    % Mark is in angle FGC. 
    \draw ($(G)!3mm!(C)$) -- ($(G)!3mm!(C) + ($(G)!3mm!(F)$) - (G)$) -- ($(G)!3mm!(F)$);
    
    % Dots for specific points
    \fill (O) circle (1.5pt);
    \fill (E) circle (1.5pt);
    \fill (F) circle (1.5pt);
    
    % Labels
    \node[above] at (A) {$A$};
    \node[left] at (B) {$B$};
    \node[right] at (C) {$C$};
    \node[above left] at (D) {$D$};
    \node[above left] at (E) {$E$};
    \node[above right] at (F) {$F$};
    \node[below right] at (G) {$G$};
    \node[below] at (H) {$H$};
    \node[below] at (O) {$O$};
\end{tikzpicture}
\end{document}
